%% LyX 2.1.4 created this file.  For more info, see http://www.lyx.org/.
%% Do not edit unless you really know what you are doing.
\documentclass[oneside,spanish]{extbook}
\usepackage[T1]{fontenc}
\usepackage[utf8]{inputenc}
\setcounter{secnumdepth}{3}
\setcounter{tocdepth}{3}
\usepackage{babel}
\addto\shorthandsspanish{\spanishdeactivate{~<>.}}
\usepackage{amsfonts}
%%\usepackage[pagebackref,ocgcolorlinks,pdfborderstyle=U]{hyperref}
%\usepackage[pagebackref,ocgcolorlinks]{hyperref}
\usepackage{graphicx}
\usepackage{commath}
\usepackage{biblatex}
\bibliography{EDA-biblio}

\begin{document}

\title{Cálculo Diferencial}

\author{Oscar Rendón Aldaraca}

\maketitle

\tableofcontents
%\listoffigures
%\listoftables

\chapter{Números reales}



\section{Los Números Reales}
\section{Axiomas de los números reales.}
\section{Intervalos y su representación gráfica.}
\section{Valor absoluto y sus propiedades.}
\section{Propiedades de las desigualdades.}
\section{Resolución de desigualdades de primer y segundo grado con una incógnita.}
\section{Resolución de desigualdades que incluyan valor absoluto.}
\chapter{Funciones.}
\section{Definición de variable, función, dominio y rango.}
\section{Función real de variable real y su representación gráfica.}
\section{Función inyectiva, suprayectiva y biyectiva.}
\section{Funciones algebraicas: polinomiales y racionales.}
\section{Funciones trascendentes: trigonométricas, logarítmicas y exponenciales.}
\section{Funciones escalonadas.}
\section{Operaciones con funciones: adición, multiplicación, división y composición.}
\section{Función inversa.}
\section{Función implícita.}
\section{Otro tipo de funciones.}

Límites y continuidad.
Noción de límite.
Definición de límite de una función.
Propiedades de los límites.
Cálculo de límites.
Límites laterales.
Límites infinitos y límites al infinito.
Asíntotas.
Continuidad en un punto y en un intervalo.
Tipos de discontinuidades.
Derivadas
Interpretación geométrica de la derivada.
Incremento y razón de cambio.
Definición de la derivada de una función.
Diferenciales.
Cálculo de derivadas.
Regla de la cadena.
Derivada de funciones implícitas.
Derivadas de orden superior.
Aplicaciones de la derivada.
Recta tangente y recta normal a una curva en un punto.
Teorema de Rolle y teoremas del valor medio.
Función creciente y decreciente.
Máximos y mínimos de una función.
Criterio de la primera derivada para máximos y mínimos.
Concavidades y puntos de inflexión.
Criterio de la segunda derivada para máximos y mínimos.
Análisis de la variación de una función. Graficación.
Problemas de optimización y de tasas relacionadas.
Cálculo de aproximaciones usando diferenciales.
La regla de L’Hôpital.



%\bibliographystyle{}



%\addbibresource{EDA-biblio}
%\bibliographystyle{plain}

\printbibliography


%\bibliography{

\end{document}

